\documentclass[11pt]{article}

\begin{document}
	
	\begin{center}
		{\Large \textbf{Online Algebra, Logic and Topology Seminar} (CMUC, University of Coimbra)}
	\end{center}
	
	\vspace{0.5cm}
	
	\noindent
	\textbf{Speaker:} Soichiro Fujii (National Institute of Informatics, Japan)
	
	\vspace{0.2cm}
	
	\noindent
	\textbf{Date:} May 4th, 2026
	
	\noindent
	\textbf{Time:} 10:00 (Europe/Lisbon)
	
	\noindent
	\textbf{Place:} Google Meet
	
	\vspace{0.5cm}
	
	\begin{center}
		{\Large \textbf{Nerves of T-categories and powerful T-functors}}
	\end{center}
	
	\vspace{0.5cm}
	
	\noindent
	\textbf{Abstract.}
	Burroni’s $T$-categories generalize internal categories, which themselves generalize ordinary (small) categories. The classical nerve construction assigns to a small category a simplicial set, and this construction extends routinely to internal categories: any internal category in a category $\mathcal{E}$ gives rise to a simplicial object in $\mathcal{E}$ as its nerve.
	
	In recent joint work with Steve Lack, titled \emph{Nerves of generalized multicategories}, we extend the nerve construction further to Burroni’s $T$-categories. For a category $\mathcal{E}$ equipped with a monad $T$, we introduce the notion of a $T$-simplicial object in $\mathcal{E}$ and show that every $T$-category determines a $T$-simplicial object. Moreover, those $T$-simplicial objects arising from $T$-categories admit a simple characterization.
	
	In this talk, I will explain these results and discuss an application to the study of powerful (or exponentiable) $T$-functors.
	
	\vspace{0.3cm}
	
	\noindent
	\textit{This talk is based on ongoing joint work with Steve Lack.}
	
\end{document}